% DRAFT METHODS ADDENDUM for the journal RLMT paper.
% This is not a standalone TeX document. Include only after the existing
% paper's RLR definitions are imported, and after method implementation
% and scientific review. NO UNOBSERVED RESULTS ARE ASSERTED HERE.
%
% Naming convention for current algorithms:
% RLR = frozen baseline; RLR-ALIGNED = new CV protocol (not presumed
% bit-identical); RLMT-R/H1 = aligned residual-screen tree;
% RLMT-LA = PROPOSED bounded child-model look-ahead (NOT implemented);
% RLMT-H2 = fixed ridge-offset ablation (not independent RLR leaves).
%
% TODO BEFORE SUBMISSION: journal-compliant bibliographic citations
% to original TSPDT, 2008 Weight-kTSP, newer weighted pair methods,
% RLR primary model, and biomedical proper-score evaluation.

\subsection{Relational logistic model tree formulation}
\label{sec:rlmt_method}

We extend the globally fitted relative logistic regression (RLR)
classifier with conditional partitions defined by within-sample
relational tests.  An internal node uses an ordered gene-pair
predicate
\begin{equation}
  z_{ab}(x)=\mathbb{I}\{x_a>x_b\},
  \qquad a\neq b ,
  \label{eq:rlmt_route}
\end{equation}
where $x$ denotes the expression vector of one sample. Unlike
a numerical threshold on an individual gene, the routing decision
depends only on the ordering of the two measured genes.  In the
present binary-only implementation the tree is deliberately
restricted to one root split and at most two terminal leaves.
Consequently, all predictions can be traced to at most one
routing relation followed by a small fitted relational model.

Each leaf $\ell\in\{L,R\}$ has its own shortlist
$\mathcal P_{\ell}$ of gene-pair relations and a penalized
logistic predictor
\begin{equation}
  p_\ell(x)=\sigma\left(
    \beta_{0,\ell}+
    \sum_{(i,j)\in\mathcal P_\ell}
      \beta_{ij,\ell}r_{ij}(x)\right),
  \label{eq:rlmt_leaf}
\end{equation}
where $r_{ij}(x)\in\{-1,0,+1\}$ retains the original RLR
margin-aware rank encoding.  The tree probability with a
validated split $z_{ab}$ is
\begin{equation}
  p_{\rm RLMT}(x)=
  z_{ab}(x)\,p_L(x)
  +\{1-z_{ab}(x)\}\,p_R(x).
  \label{eq:rlmt_predict}
\end{equation}
The degenerate no-split model uses one global parent predictor
$p_0(x)$. The tree formulation therefore generalizes the
single-model representation structurally, but numerical identity
with the frozen original RLR baseline requires an identical
training algorithm and validation protocol.

\subsection{Training-only split discovery and confirmation}
\label{sec:rlmt_discovery}

To prevent candidate selection from exploiting the outer test
samples, all supervised operations are nested inside the
corresponding outer training set $T$.  We partition
$T=D\cup V$ into a discovery subset $D$ and disjoint
confirmation subset $V$. A common, deterministic candidate
set $\mathcal C_D$ of relational predicates is generated from
$D$ only, subject to both minimum child size and minimum
class-specific child support. Class balance is evaluated from
labels in $D$ only. For a candidate $c\in\mathcal C_D$,
let $\widehat{p}_{D,c}(x)$ denote the two-leaf predictor fitted
on the discovery data and $\widehat{p}_{D,0}(x)$ the parent
model fitted by the same training procedure.

We distinguish three candidate-scoring strategies for a
controlled methodological comparison:
\begin{enumerate}
  \item \textit{TSPDT-style baseline}: select an admissible
  relational predicate by a conventional impurity reduction
  evaluated in $D$, then fit the same local RLR learner.
  \item \textit{Residual-guided RLMT-R}: identify a
  promising routing relation using the association between
  class labels, routed groups, and parent out-of-fold
  residuals computed exclusively within $D$.
  \item \textit{Look-ahead RLMT-LA}: use a bounded shortlist
  of candidates scored in $D$ by the predicted gain in
  probabilistic loss after actually fitting the prospective
  RLR leaf models, rather than optimizing impurity alone.
  The exact shortlist budget must be prespecified on
  synthetic-only tests before biomedical outcomes are
  examined. \textbf{Implementation pending.}
\end{enumerate}

A single nominated candidate is evaluated on $V$. With the
binary logloss
\begin{equation}
  \mathcal L(y,p)=-y\log p-(1-y)\log(1-p)
  \label{eq:rlmt_logloss}
\end{equation}
and paired per-sample confirmation improvement
\begin{equation}
  \Delta_i=
  \mathcal L(y_i,\widehat p_{D,0}(x_i))
  -
  \mathcal L(y_i,\widehat p_{D,c}(x_i)),
  \quad i\in V ,
  \label{eq:rlmt_paired_delta}
\end{equation}
the split is accepted only when the prespecified conservative
confirmation lower bound exceeds zero and adequate validation
support is available.  Candidate search is not repeated after
inspecting $V$, and there is no outer-test-based pruning.
If confirmation rejects the split, the final predictor is the
one-leaf parent fitted on the full outer training set $T$.

\subsection{Consistent discovery and final estimators}
\label{sec:rlmt_alignment}

An essential implementation requirement is that each learning
stage fit exactly the same model \emph{procedure} when applied
to a different training subset. In particular, define the
parent-fitting operator
\begin{equation}
 A(S): S\mapsto\{
   \mathcal P_S,\widehat\beta_S,
   \widehat\lambda_S,\widehat p_S\},
 \label{eq:rlmt_fit_operator}
\end{equation}
which performs supervised pair ranking, gene-disjoint shortlist
construction, rank-margin encoding, and elastic-net
regularization. The same operator and deterministic inner-CV
rule are applied to $S=D$ for confirmation and to $S=T$ for
final training, with all supervised pair-selection steps
recomputed using the supplied subset $S$ only.
Thus, the \emph{fitting procedure} is invariant, although
the fitted relations and coefficients can of course differ
between discovery and full outer training sets.

The experimental aligned parent, denoted
\textsc{RLR-Aligned}, uses explicit stratified inner
five-fold CV and is kept separate from the exact,
archived RLR comparator until prediction-level parity
has been demonstrated. Reported comparisons must include
both controls to prevent inadvertently crediting changes
in the root fitting procedure to tree structure.

\subsection{Offset-corrected leaf ablation}
\label{sec:rlmt_offset}

As a separate ablation we evaluate conservative
leaf-specific corrections to the aligned global model.
For leaf $\ell$, an offset model estimates
\begin{equation}
  p^{\rm off}_{\ell}(x)=
  \sigma\left[
    \operatorname{logit}\{\widehat p_0(x)\}
    +\delta_{\ell}(x)
  \right],
  \label{eq:rlmt_offset}
\end{equation}
where $\delta_{\ell}(x)$ is a fixed-penalty ridge regression
on selected relational features, with no freely estimated
unpenalized intercept. Training offsets use \emph{honest
out-of-fold parent probabilities} within the leaf-training
data; they must not use the parent's in-sample fitted scores
as if those were independently predicted outcomes. The
current predeclared research configuration uses
$\alpha=0$, $\lambda=1$, and no within-leaf lambda
cross-validation; this is a fixed design parameter rather
than a selection made from old held-out biomedical losses.

This approach constitutes a \emph{different model class}
from independent RLR leaf models in
Eq.~\eqref{eq:rlmt_leaf}. It must therefore appear as
an explicitly named offset ablation, not as a replacement
for the primary relational model-tree contribution.

\subsection{Reproducibility and evidence boundaries}
\label{sec:rlmt_evidence}

For every outer split we preserve the fitted parent
and, when accepted, both complete leaf models; the selected
routing relation and confirmation decision; all class
support counts; the training-only inner fold identities;
and sample-aligned predicted probabilities. The archived
original RLR results remain unchanged. We distinguish
retrospective development analyses on previously exposed
microarray cohorts from later independently selected
clinical-cohort evaluation. The primary proper scoring
rule is unweighted outer-test logloss, complemented by
Brier score, class-conditional errors, AUC, and
thresholded classification measures. No claim of superior
predictive performance is made unless supported by the
completed paired evaluations.
